\section{Benchmark 体系升级：从网页基准到 OS-World}

\subsection{现有基准的价值与边界}
讲者回顾了 Mind2Web、WebArena、VisualWebArena 等基准。这些工作推动了 Web Agent 研究，但它们的环境覆盖、任务类型和可扩展性仍不足以代表真实 ``computer use'' 场景，尤其是跨应用与开放式目标任务。

\begin{knowledgebox}{典型不足}
\begin{itemize}
    \item 任务分布偏网页导航，覆盖面窄；
    \item 许多任务存在隐式模板化路径，容易被策略过拟合；
    \item 对桌面应用、文件系统、跨软件流程支持不足；
    \item 难以承载 ``开发-测试-评测'' 一体化闭环。
\end{itemize}
\end{knowledgebox}

\subsection{OS-World 的设计目标}
OS-World 被定位为更通用的真实计算机环境：在虚拟机中提供可交互桌面与应用生态，支持研究者定义任务、执行 Agent、收集轨迹、运行评测脚本，并且保留扩展空间。

\begin{importantbox}{OS-World 的三项关键属性}
\begin{enumerate}
    \item \textbf{真实性}：接近真实操作系统与应用行为；
    \item \textbf{可扩展性}：可持续新增任务与应用组合；
    \item \textbf{可评测性}：每个任务可绑定脚本化验证逻辑。
\end{enumerate}
\end{importantbox}

\begin{figure}[H]
\centering
\includegraphics[width=0.88\textwidth]{frame-02-osworld.jpg}
\caption{视频约 00:10:32：OS-World 作为可扩展真实环境被提出，用于多模态 Agent 的开发与评测}
\end{figure}

\subsection{对比视角：为什么 OS-World 更接近生产问题}
\begin{table}[H]
\centering
\begin{tabular}{p{3.2cm}p{3.7cm}p{3.8cm}}
\toprule
\textbf{维度} & \textbf{传统 Web 基准} & \textbf{OS-World} \\
\midrule
环境覆盖 & 主要网页域 & 网页 + 桌面应用 + 系统交互 \\
任务形态 & 偏导航与表单 & 跨应用、多步骤、开放式流程 \\
评测方式 & 页面结果导向 & 脚本验证 + 部分完成奖励 \\
扩展能力 & 新任务构建成本较高 & 支持系统化新增任务与配置 \\
工程适配 & 偏离线评测 & 开发、训练、评测一体化 \\
\bottomrule
\end{tabular}
\caption{OS-World 并非替代所有基准，而是补齐 ``真实计算机任务'' 的核心空缺}
\end{table}

\begin{warningbox}{评测迁移风险}
如果团队只在单一 benchmark 上优化，模型容易对评测脚本形成 ``策略投机''。课程建议在多环境、多任务族上做交叉验证，以防止 benchmark overfitting。
\end{warningbox}

\subsection{本章小结}
OS-World 的意义不只是多一个 benchmark，而是把 Agent 研究从 ``题库竞争'' 推向 ``环境工程''。这为后续的数据与模型设计提供了统一实验底座。

